% ============================================================
%  SY02-28 《p 进局部系统的刚性与 Riemann–Hilbert 对应》译稿导言区
%  底本：Ruochan Liu, Xinwen Zhu, Rigidity and a Riemann–Hilbert
%        correspondence for p-adic local systems, arXiv:1602.06282v3
%        （35 页；Inventiones mathematicae 209 (2017), 291–333）
%  体例：定理与注各自按节编号（Theorem 1.1、Remark 1.1 并行）；
%        公式按节编号（2.1）。
% ============================================================

% ---------- 页面与字体 ----------
\usepackage{geometry}
\geometry{a4paper, left=2.6cm, right=2.6cm, top=2.8cm, bottom=2.6cm}
\linespread{1.25}
\setlength{\parskip}{0.28em}
\setlength{\emergencystretch}{1.5em}
\sloppy

% ---------- 数学 ----------
\usepackage{amsmath}
\usepackage{amsthm}
\usepackage{amssymb}
\usepackage{amsfonts}
\usepackage{mathtools}
\usepackage{mathrsfs}
\usepackage{stmaryrd}
\usepackage{bm}
\allowdisplaybreaks[2]

% ---------- 图表与列表 ----------
\usepackage{graphicx}
\usepackage{xcolor}
\usepackage{enumitem}
\setlist{itemsep=2pt, topsep=3pt, parsep=0pt}
\setlist[enumerate,1]{label=(\arabic*)}

% ---------- 超链接 ----------
\usepackage[hyperfootnotes=false, colorlinks=true, linkcolor=blue!45!black,
            citecolor=green!35!black, urlcolor=teal!60!black,
            bookmarksnumbered=true]{hyperref}
\usepackage{bookmark}

% ---------- 中文微调 ----------
\xeCJKsetup{CJKecglue={\hskip 0.24em plus 0.05em minus 0.05em}}
\ctexset{
  section    = {name={,}, aftername={.\enspace}, format=\Large\bfseries\centering},
  subsection = {aftername={.\enspace}, format=\large\bfseries},
  subsubsection = {aftername={.\enspace}, format=\normalsize\bfseries}
}

% ---------- 定理环境：定理与注各自按节编号，与原文一致 ----------
\numberwithin{equation}{section}
\theoremstyle{plain}
\newtheorem{thm}{定理}[section]
\newtheorem{prop}{命题}[section]
\newtheorem{lem}{引理}[section]
\newtheorem{cor}{推论}[section]
\newtheorem*{conj*}{猜想}
\theoremstyle{definition}
\newtheorem{dfn}{定义}[section]
\newtheorem{construction}{构造}[section]
\theoremstyle{remark}
\newtheorem{rmk}{注}[section]
\newtheorem{ex}{例}[section]

\renewcommand{\proofname}{证明}
\renewcommand{\qedsymbol}{$\square$}
\renewcommand{\refname}{参考文献}
\renewcommand{\contentsname}{目\quad 录}

% ---------- 译注标记（独立脚注序列） ----------
\usepackage{manyfoot}
\DeclareNewFootnote[plain]{trans}
\renewcommand{\thefootnotetrans}{译\arabic{footnotetrans}}
\newcommand{\zy}[1]{\footnotetrans{#1}}

% ---------- 常用记号 ----------
\newcommand{\Q}{\mathbb{Q}}
\newcommand{\Z}{\mathbb{Z}}
\newcommand{\R}{\mathbb{R}}
\newcommand{\C}{\mathbb{C}}
\newcommand{\Qp}{\mathbb{Q}_p}
\newcommand{\Zp}{\mathbb{Z}_p}
\newcommand{\Ql}{\mathbb{Q}_\ell}
\newcommand{\Zl}{\mathbb{Z}_\ell}
\newcommand{\Af}{\mathbb{A}_f}
\newcommand{\Gal}{\operatorname{Gal}}
\newcommand{\dR}{\mathrm{dR}}
\newcommand{\BdR}{B_{\mathrm{dR}}}
\newcommand{\BdRp}{B_{\mathrm{dR}}^{+}}
\newcommand{\OX}{\mathcal{O}_X}
\newcommand{\OXB}{\mathcal{O}_X\widehat{\otimes}\BdR}
\newcommand{\Xk}{\widehat{X}_{\bar k}}
\newcommand{\kb}{\bar k}
\newcommand{\khb}{\widehat{\bar k}}
\newcommand{\proet}{\text{pro\'et}}
\newcommand{\Fil}{\mathrm{Fil}}
\newcommand{\gr}{\mathrm{gr}}
\newcommand{\PerDR}{P_{\dR}}
\newcommand{\Sh}{\operatorname{Sh}}
\newcommand{\Spa}{\operatorname{Spa}}
\newcommand{\Spec}{\operatorname{Spec}}
\newcommand{\id}{\mathrm{id}}
\newcommand{\on}[1]{\operatorname{#1}}
\newcommand{\sgn}{\operatorname{sgn}}
\newcommand{\Tr}{\operatorname{Tr}}

% 引用键统一走 \cite；参考文献条目保留原文
\providecommand{\etal}{\textit{et al.}}
